\newcommand{\figuraMareasVivasMuertas}{%
\begin{figure}[H]
    \centering
    \begin{tikzpicture}[>=Latex,font=\small,line cap=round,line join=round]
        \coordinate (S1) at (1.2,4.8);
        \coordinate (T1) at (7.0,4.8);
        \coordinate (S2) at (1.2,1.7);
        \coordinate (T2) at (7.0,1.7);

        \node[font=\bfseries] at (6.2,6.15)
            {(a) Mareas vivas: alineación};
        \node[font=\bfseries] at (6.2,3.22)
            {(b) Mareas muertas: cuadratura};

        \foreach \S in {S1,S2} {
            \fill[yellow!65] (\S) circle (0.60);
            \draw[very thick,orange!85!black] (\S) circle (0.60);
        }
        \foreach \T in {T1,T2} {
            \fill[blue!25] (\T) circle (0.48);
            \draw[very thick,blue!70!black] (\T) circle (0.48);
        }

        \draw[very thick,red!70!black] (T1) ellipse (1.12 and 0.66);
        \draw[very thick,red!70!black] (T2) ellipse (0.83 and 0.65);

        \node at (1.2,3.95) {Sol};
        \node at (1.2,0.85) {Sol};
        \node at (7.0,3.88) {Tierra};
        \node[anchor=east] at (6.35,1.25) {Tierra};

        \draw[very thick,orange!70!black] (1.85,4.8)--(10.65,4.8);
        \fill[gray!45] (5.65,4.8) circle (0.20);
        \draw[thick] (5.65,4.8) circle (0.20);
        \node[anchor=south] at (5.65,5.10) {Luna nueva};
        \fill[gray!45] (8.35,4.8) circle (0.20);
        \draw[thick] (8.35,4.8) circle (0.20);
        \node[anchor=south] at (8.35,5.10) {Luna llena};
        \draw[<->,red!75!black,ultra thick] (6.08,4.8)--(7.92,4.8);
        \node[red!75!black,anchor=west] at (8.55,4.35)
            {rango máximo};

        \draw[very thick,orange!70!black] (1.85,1.7)--(10.65,1.7);
        \draw[very thick,purple!70!black] (7.0,0.35)--(7.0,2.85);
        \fill[gray!45] (7.0,2.55) circle (0.20);
        \draw[thick] (7.0,2.55) circle (0.20);
        \fill[gray!45] (7.0,0.85) circle (0.20);
        \draw[thick] (7.0,0.85) circle (0.20);
        \node[anchor=west] at (7.75,2.62) {cuarto creciente};
        \node[anchor=west] at (7.75,0.72) {cuarto menguante};
        \draw[<->,red!75!black,thick] (6.20,1.7)--(7.80,1.7);
        \draw[<->,purple!70!black,thick] (7.0,1.02)--(7.0,2.38);
        \node[align=center,text=black!65] at (3.75,0.55)
            {los efectos solar y lunar\\se compensan parcialmente};
    \end{tikzpicture}
    \caption{Superposición idealizada de las mareas lunar y solar. En la
    alineación, ambas deformaciones tienen el mismo eje y se refuerzan; en la
    cuadratura, sus ejes son perpendiculares y el rango de marea disminuye.
    Las distancias y los tamaños no están a escala.}
    \label{fig:mareas-vivas-muertas}
\end{figure}%
}